% Source: https://artofproblemsolving.com/community/c3h1447778_diagonals
% Stable identity: aops-topic-1447778-q01
\textbf{(a)} Label the vertices $P_1,\dots,P_n$. From a vertex, a diagonal can go to any of the other $n-1$ vertices except its two neighbours, so there are $n-3$ diagonals at each vertex. For $n=7$ the $4$ diagonals from $P_1$ are shown in red:

\begin{asy}
usepackage("amsmath");
size(4cm);
defaultpen(fontsize(10pt));
pair[] P={(0.0000,2.0000),(-1.5637,1.2470),(-1.9499,-0.4450),(-0.8678,-1.8019),(0.8678,-1.8019),(1.9499,-0.4450),(1.5637,1.2470)};
draw(P[0]--P[2],red);
draw(P[0]--P[3],red);
draw(P[0]--P[4],red);
draw(P[0]--P[5],red);
draw(P[0]--P[1]--P[2]--P[3]--P[4]--P[5]--P[6]--cycle,linewidth(1));
for(int i=0;i<7;++i) dot(P[i]);
label("$P_{1}$",P[0],dir(P[0]));
label("$P_{2}$",P[1],dir(P[1]));
label("$P_{3}$",P[2],dir(P[2]));
label("$P_{4}$",P[3],dir(P[3]));
label("$P_{5}$",P[4],dir(P[4]));
label("$P_{6}$",P[5],dir(P[5]));
label("$P_{7}$",P[6],dir(P[6]));
label("\mbox{}",(min(currentpicture,user=true)+max(currentpicture,user=true))/2,black); // keeps the text after the figure black when LaTeX typesets the labels inline
\end{asy}

Adding over all $n$ vertices gives $n(n-3)$, but each diagonal has two ends and has been counted twice. Hence the number of diagonals is
\[
\boxed{\dfrac{n(n-3)}{2}} .
\]
(Equivalently: $\dbinom{n}{2}$ segments join two vertices, and $n$ of them are sides, so there are $\dbinom{n}{2}-n=\dfrac{n(n-3)}{2}$ diagonals.)

\textbf{(b)} For $n=1000$:
\[
\dfrac{1000\cdot997}{2}=\boxed{498500}.
\]
